% Options for packages loaded elsewhere
\PassOptionsToPackage{unicode}{hyperref}
\PassOptionsToPackage{hyphens}{url}
%
\documentclass[
  11pt,
]{article}
\usepackage{amsmath,amssymb}
\usepackage{iftex}
\ifPDFTeX
  \usepackage[T1]{fontenc}
  \usepackage[utf8]{inputenc}
  \usepackage{textcomp} % provide euro and other symbols
\else % if luatex or xetex
  \usepackage{unicode-math} % this also loads fontspec
  \defaultfontfeatures{Scale=MatchLowercase}
  \defaultfontfeatures[\rmfamily]{Ligatures=TeX,Scale=1}
\fi
\usepackage{lmodern}
\ifPDFTeX\else
  % xetex/luatex font selection
  \setmainfont[]{Calibri}
\fi
% Use upquote if available, for straight quotes in verbatim environments
\IfFileExists{upquote.sty}{\usepackage{upquote}}{}
\IfFileExists{microtype.sty}{% use microtype if available
  \usepackage[]{microtype}
  \UseMicrotypeSet[protrusion]{basicmath} % disable protrusion for tt fonts
}{}
\makeatletter
\@ifundefined{KOMAClassName}{% if non-KOMA class
  \IfFileExists{parskip.sty}{%
    \usepackage{parskip}
  }{% else
    \setlength{\parindent}{0pt}
    \setlength{\parskip}{6pt plus 2pt minus 1pt}}
}{% if KOMA class
  \KOMAoptions{parskip=half}}
\makeatother
\usepackage{xcolor}
\usepackage[margin=1in]{geometry}
\usepackage{longtable,booktabs,array}
\usepackage{calc} % for calculating minipage widths
% Correct order of tables after \paragraph or \subparagraph
\usepackage{etoolbox}
\makeatletter
\patchcmd\longtable{\par}{\if@noskipsec\mbox{}\fi\par}{}{}
\makeatother
% Allow footnotes in longtable head/foot
\IfFileExists{footnotehyper.sty}{\usepackage{footnotehyper}}{\usepackage{footnote}}
\makesavenoteenv{longtable}
\usepackage{graphicx}
\makeatletter
\def\maxwidth{\ifdim\Gin@nat@width>\linewidth\linewidth\else\Gin@nat@width\fi}
\def\maxheight{\ifdim\Gin@nat@height>\textheight\textheight\else\Gin@nat@height\fi}
\makeatother
% Scale images if necessary, so that they will not overflow the page
% margins by default, and it is still possible to overwrite the defaults
% using explicit options in \includegraphics[width, height, ...]{}
\setkeys{Gin}{width=\maxwidth,height=\maxheight,keepaspectratio}
% Set default figure placement to htbp
\makeatletter
\def\fps@figure{htbp}
\makeatother
\setlength{\emergencystretch}{3em} % prevent overfull lines
\providecommand{\tightlist}{%
  \setlength{\itemsep}{0pt}\setlength{\parskip}{0pt}}
\setcounter{secnumdepth}{-\maxdimen} % remove section numbering
\ifLuaTeX
  \usepackage{selnolig}  % disable illegal ligatures
\fi
\IfFileExists{bookmark.sty}{\usepackage{bookmark}}{\usepackage{hyperref}}
\IfFileExists{xurl.sty}{\usepackage{xurl}}{} % add URL line breaks if available
\urlstyle{same}
\hypersetup{
  pdftitle={Recamán, Collatz, and Delannoy: An Obstruction-Bit Taxonomy},
  pdfauthor={Zolton (synthesis from collaborative drafts)},
  hidelinks,
  pdfcreator={LaTeX via pandoc}}

\title{Recamán, Collatz, and Delannoy: An Obstruction-Bit Taxonomy}
\usepackage{etoolbox}
\makeatletter
\providecommand{\subtitle}[1]{% add subtitle to \maketitle
  \apptocmd{\@title}{\par {\large #1 \par}}{}{}
}
\makeatother
\subtitle{Markov Shadows, Operator Wheels, and the 321210 Pattern}
\author{Zolton (synthesis from collaborative drafts)}
\date{May 2026}

\begin{document}
\maketitle

{
\setcounter{tocdepth}{2}
\tableofcontents
}
\hypertarget{abstract}{%
\section{Abstract}\label{abstract}}

We propose a taxonomy of integer-sequence dynamical systems by
\emph{where the obstruction information lives}: locally in the current
state (Collatz), globally in the entire history (Recamán), or freely as
an unconstrained choice (Delannoy walks). All three admit a
Collatz-style branch-map form \[
x_{n+1} = F_{b_n}(x_n), \qquad b_n \in \{0,1\},
\] and the differences in how \(b_n\) is generated explain the
qualitative differences in their dynamics.

For Recamán specifically we tested whether a finite-state symbolic
shadow of the \textbf{obstruction process} \(\{b_n\}\) using the
reverse-complement involution \(\Theta_3\) on words over \(\{0,1,2,3\}\)
captures the dynamics. The empirical answer is \textbf{no}: the
two-state \(\Theta_3\) wheel has essentially zero predictive content for
\(b_n\) (\(|q(0) - q(1)| < 5 \times 10^{-4}\)), and KL divergence on
\(8\)-block bit distributions exceeds the Markov-null noise floor by a
factor of \(\sim 3{,}000\). What does work is bit-history conditioning:
\(\Pr[b_n \ne b_{n-1}] = 0.993\), i.e.~the obstruction stream is
\textbf{near-perfect alternation with rare phase-slip defects at rate
\(\sim 7 \times 10^{-3}\)}. The remaining task is therefore to
characterise the phase-slip process, not to fit a Markov wheel.

We formulate the \textbf{Recamán Markov-Shadow Conjecture} on the value
side: late or missing values are statistically enriched near an
approximate Markov surface \[
X^2 + Y^2 + Z^2 \;\approx\; \kappa\, XYZ.
\] A \(\zeta(3)\)-coupling is shown to be impossible over integers by
Apéry's theorem; only an approximate real surface is admissible.
Delannoy magnitudes appear in the Recamán positional distribution
\textbf{because the Recamán obstruction process generates them}: the
three-way branch structure (down-free / down-blocked-becomes-up /
up-by-boundary) is intrinsic to the rule, and the Delannoy recurrence
falls out of the operator-word counting. We end with a falsifiable
empirical protocol and open problems.

\begin{quote}
\textbf{Status.} The taxonomy in Section 2 is a genuine conceptual
claim. The \(\Theta_3\) wheel has been \textbf{empirically falsified} as
a Markov model of the obstruction process; the replacement is
bit-history conditioning and a sharpened open problem on phase-slip
dynamics (Section 3.5). The Markov-shadow conjecture on the value side
is empirically testable but currently untested on the full Chan hole
list. The Delannoy structure in the positional distribution is generated
by Recamán's own three-way branch and is \emph{internal} to the
dynamics, not an external envelope.
\end{quote}

\hypertarget{setup-three-sequences-one-branch-map-form}{%
\section{1. Setup: Three Sequences, One Branch-Map
Form}\label{setup-three-sequences-one-branch-map-form}}

\hypertarget{recamuxe1n}{%
\subsection{1.1 Recamán}\label{recamuxe1n}}

Recamán's sequence is defined by \(a_0 = 0\) and \[
a_n = \begin{cases} a_{n-1} - n & \text{if } a_{n-1} - n > 0 \text{ and } a_{n-1} - n \notin V_{n-1}, \\ a_{n-1} + n & \text{otherwise,} \end{cases}
\] where \(V_{n-1} = \{a_0, a_1, \dots, a_{n-1}\}\).

Introducing the \textbf{obstruction bit} \[
b_n \;=\; \mathbf{1}\!\left[\,a_{n-1} - n \le 0 \;\;\lor\;\; a_{n-1} - n \in V_{n-1}\,\right],
\] the recursion collapses to the compact Collatz-style form \[
\boxed{\;a_n \;=\; a_{n-1} + n\,(2 b_n - 1)\;}
\]

\hypertarget{the-signed-triangular-identity-and-the-rigorous-tower}{%
\subsubsection{1.1.1 The signed-triangular identity and the rigorous
tower}\label{the-signed-triangular-identity-and-the-rigorous-tower}}

Unrolling the recursion gives the \textbf{exact closed form} \[
\boxed{\;a_n \;=\; T_n \;-\; 2 \sum_{i \in D_n} i, \qquad T_n = \frac{n(n+1)}{2},\;}
\] where \(D_n \subseteq \{1, \dots, n\}\) is the set of indices at
which the rule took a \emph{down} step (equivalently,
\(i \in D_n \iff b_i = 0\)). This identity is exact, not heuristic, and
we have verified it numerically up to \(n = 2 \times 10^5\).

It supplies the \textbf{rigorous, non-numerological version of the tower
idea}: rather than decomposing \(a_n\) by its decimal digits, we
decompose it by the deterministic objects that produced it. The layers
are \[
\begin{aligned}
L_0 &= a_n & &\text{(the value)} \\
L_1 &= T_n = \tfrac{n(n+1)}{2} & &\text{(triangular envelope)} \\
L_2 &= D_n & &\text{(down-step index set)} \\
L_3 &= a_{n-1} - n & &\text{(candidate down-target)} \\
L_4 &= \mathbf{1}[a_{n-1} - n \in V_{n-1}] & &\text{(collision indicator)} \\
L_5 &= \text{positional encoding / residue wheel} & &\text{(symbolic shadow)}
\end{aligned}
\] Each layer is generated from \((a_{n-1}, V_{n-1}, n)\) by a single
arithmetic or logical operation; the layers are stacked so that
\emph{every dynamical question about Recamán reduces to a question about
how \(D_n\) is built}. The Markov-shadow defect of Section 4, the
obstruction probabilities \(q(w)\) of Section 6, and the Delannoy
generation of Section 7 are all statements about the structure of
\(D_n\) in disguise.

\hypertarget{collatz}{%
\subsection{1.2 Collatz}\label{collatz}}

The Collatz map is \[
T(x) = \begin{cases} x/2 & x \equiv 0 \pmod 2, \\ (3x+1)/2 & x \equiv 1 \pmod 2. \end{cases}
\] Setting \(b_n = x_n \bmod 2\), \[
x_{n+1} \;=\; b_n \cdot \frac{3x_n + 1}{2} \;+\; (1 - b_n) \cdot \frac{x_n}{2}.
\]

\hypertarget{delannoy-walks}{%
\subsection{1.3 Delannoy walks}\label{delannoy-walks}}

A Delannoy walk on \(\mathbb{Z}_{\ge 0}^2\) from \((0,0)\) to \((m,n)\)
takes steps in \(\{E, N, NE\}\). The number of such walks is the
Delannoy number \(D(m,n)\), satisfying \[
D(m,n) \;=\; D(m-1, n) + D(m, n-1) + D(m-1, n-1).
\] Here the ``branch bit'' generalises to a ternary choice
\(c_n \in \{E, N, NE\}\), freely selected.

\hypertarget{the-obstruction-bit-taxonomy}{%
\section{2. The Obstruction-Bit
Taxonomy}\label{the-obstruction-bit-taxonomy}}

The three systems are unified --- and distinguished --- by where the
branch information lives.

\begin{longtable}[]{@{}
  >{\raggedright\arraybackslash}p{(\columnwidth - 6\tabcolsep) * \real{0.11}}
  >{\raggedright\arraybackslash}p{(\columnwidth - 6\tabcolsep) * \real{0.20}}
  >{\raggedright\arraybackslash}p{(\columnwidth - 6\tabcolsep) * \real{0.43}}
  >{\raggedright\arraybackslash}p{(\columnwidth - 6\tabcolsep) * \real{0.26}}@{}}
\toprule\noalign{}
\begin{minipage}[b]{\linewidth}\raggedright
System
\end{minipage} & \begin{minipage}[b]{\linewidth}\raggedright
State
\end{minipage} & \begin{minipage}[b]{\linewidth}\raggedright
Obstruction bit / choice
\end{minipage} & \begin{minipage}[b]{\linewidth}\raggedright
Information locality
\end{minipage} \\
\midrule\noalign{}
\endhead
\bottomrule\noalign{}
\endlastfoot
Collatz & \(x \in \mathbb{Z}_+\) & \(b_n = x_n \bmod 2\) &
\textbf{Local} --- one bit of current state \\
Recamán & \((a_{n-1}, n, V_{n-1})\) &
\(b_n = \mathbf{1}[a_{n-1}{-}n \le 0 \,\lor\, a_{n-1}{-}n \in V_{n-1}]\)
& \textbf{Global} --- entire history \\
Delannoy & \((i,j) \in \mathbb{Z}_+^2\) & \(c_n \in \{E, N, NE\}\) free
& \textbf{External} --- chosen, not derived \\
\end{longtable}

\textbf{The conceptual claim.} Many ``wild'' integer sequences are best
understood not by their values but by the \emph{information structure of
their branch rule}. Collatz is hard because the local bit interacts
nontrivially with the multiplicative dynamics. Recamán is hard because
the bit is non-local. Delannoy is tractable because there is no
constraint to satisfy --- every path is legal.

This taxonomy clarifies what \emph{kind} of theorem one can expect for
each system:

\begin{itemize}
\tightlist
\item
  \textbf{Local-bit systems} admit ergodic-theoretic and 2-adic analysis
  (the Collatz literature).
\item
  \textbf{Global-bit systems} require a closure approximation: replace
  the true non-local bit by a functional of a low-dimensional shadow.
  This is the role of the operator wheel below.
\item
  \textbf{Free-choice systems} are counted by combinatorial recurrences.
  They are not \emph{external} envelopes on constrained systems ---
  rather, the constrained systems \emph{generate} free-choice structure
  when their constraints take a three-way form.
\end{itemize}

The Recamán--Delannoy connection sits inside Recamán: the obstruction
rule at each step partitions into a three-way local branch (down-free /
down-blocked-forces-up / up-by-boundary), and the operator-word count
generated by Recamán's own dynamics satisfies a Delannoy recurrence.
Delannoy magnitudes appear in the positional distribution
\textbf{because the Recamán process produces them}, not because they
bound it from outside. We make this precise in Section 7.

\begin{quote}
\textbf{Remark.} Variants of this taxonomy appear throughout
symbolic-dynamics and ergodic-theory literature under different names
(skew products, random-environment substitutions). The novelty here is
the specific three-way comparison and using it to predict \emph{which}
analytical tools should work on Recamán.
\end{quote}

\hypertarget{the-reverse-complement-wheel-theta_3}{%
\section{\texorpdfstring{3. The Reverse-Complement Wheel
\(\Theta_3\)}{3. The Reverse-Complement Wheel \textbackslash Theta\_3}}\label{the-reverse-complement-wheel-theta_3}}

\hypertarget{the-reverse-complement-involution-theta_3}{%
\subsection{\texorpdfstring{3.1 The reverse-complement involution
\(\Theta_3\)}{3.1 The reverse-complement involution \textbackslash Theta\_3}}\label{the-reverse-complement-involution-theta_3}}

On words over \(\{0,1,2,3\}\), define the \textbf{reverse-complement
involution} \[
\Theta_3(w_1 w_2 \cdots w_m) \;=\; \operatorname{rev}\!\left((3 - w_1)(3 - w_2) \cdots (3 - w_m)\right).
\] In particular, \[
\Theta_3(210) = 321, \qquad \Theta_3(321) = 210,
\] so the six-digit word \(321210 = 321\,|\,210\) is a fixed point of
the split condition \(L = \Theta_3(R)\), since \(3+0 = 2+1 = 1+2 = 3\).

This is a clean symbolic symmetry of the alphabet \(\{0,1,2,3\}\). The
remaining question is whether it usefully describes Recamán's
obstruction process. We address this empirically.

\hypertarget{the-two-state-wheel-and-its-proposed-master-equation}{%
\subsection{3.2 The two-state wheel and its proposed master
equation}\label{the-two-state-wheel-and-its-proposed-master-equation}}

The natural two-state model uses \(S_n \in \{210, 321\}\) with the
transition \[
S_{n+1} \;=\; \begin{cases} \Theta_3(S_n) & \text{if } b_n = 1, \\ S_n & \text{if } b_n = 0. \end{cases}
\] If this were a Markov model of the obstruction process, the
conditional probabilities \(q(w) = \Pr[b_n = 1 \mid S_n = w]\) would
satisfy a master equation with stationary distribution \[
p(210) \;=\; \frac{q_{321}}{q_{210} + q_{321}}, \qquad p(321) \;=\; \frac{q_{210}}{q_{210} + q_{321}}.
\] \textbf{Whether the model is appropriate is an empirical question},
settled in 3.3.

\hypertarget{empirical-test-of-the-wheel-a-strong-negative-result}{%
\subsection{3.3 Empirical test of the wheel: a strong negative
result}\label{empirical-test-of-the-wheel-a-strong-negative-result}}

We ran the true Recamán generator to \(N = 2 \times 10^5\) terms,
extracted the real obstruction stream \(\{b_n\}\) and the wheel-state
stream \(\{S_n\}\), and measured the conditional block probabilities
directly:

\[
q(0) \;=\; 0.5000, \qquad q(1) \;=\; 0.5001, \qquad |q(0) - q(1)| \;<\; 5 \times 10^{-4}.
\]

The wheel state has \textbf{essentially zero predictive content} for the
obstruction bit. The two-state \(\Theta_3\) shadow does not separate
Recamán's obstruction events.

We confirmed this with three independent tests:

\textbf{(i) Mutual information.} \(I(b_n \,;\, S_{n-1}) = 0.0000\) bits
at both \(N = 50{,}000\) and \(N = 200{,}000\), against a marginal
entropy \(H(b) = 1.0000\) bits.

\textbf{(ii) Run-length distribution.} Real Recamán obstruction bits
produce length-1 runs at \(\sim 4\times\) the rate of a Markov chain
with matched marginals, and produce essentially zero runs of length
\(\ge 6\):

\begin{longtable}[]{@{}llll@{}}
\toprule\noalign{}
Run length & Real (N=200k) & Markov null & Ratio \\
\midrule\noalign{}
\endhead
\bottomrule\noalign{}
\endlastfoot
1 & \(197{,}699\) & \(50{,}087\) & \(3.95\times\) \\
2 & \(326\) & \(24{,}951\) & \(0.013\times\) \\
3 & \(422\) & \(12{,}547\) & \(0.034\times\) \\
\(\ge 4\) & \(93\) & \(12{,}049\) & \(0.008\times\) \\
\end{longtable}

\textbf{(iii) KL divergence on \(k\)-block distributions.} At \(k = 8\)
block length, \[
D_{\text{KL}}\bigl(\text{real} \,\|\, \text{Markov null}\bigr) \;=\; 4.63 \text{ nats},
\] versus a noise floor of \(0.0015\) nats between independent null
realisations --- an excess of \(\sim 3{,}000\times\) the baseline.

\textbf{Conclusion.} The obstruction process is provably non-Markovian
on the \(\Theta_3\) wheel. Any model that conditions \(b_n\) on the
wheel state alone cannot reproduce Recamán's bit-level structure.

\hypertarget{what-replaces-the-wheel-bit-history-conditioning}{%
\subsection{3.4 What replaces the wheel: bit-history
conditioning}\label{what-replaces-the-wheel-bit-history-conditioning}}

A held-out validation across a broad family of candidate states (we
split \(N\) in half, fit \(\Pr[b_n = 1 \mid X]\) on the first half,
scored log-likelihood on the second) identified what \emph{does} carry
predictive information:

\begin{longtable}[]{@{}
  >{\raggedright\arraybackslash}p{(\columnwidth - 4\tabcolsep) * \real{0.3333}}
  >{\raggedright\arraybackslash}p{(\columnwidth - 4\tabcolsep) * \real{0.3333}}
  >{\raggedright\arraybackslash}p{(\columnwidth - 4\tabcolsep) * \real{0.3333}}@{}}
\toprule\noalign{}
\begin{minipage}[b]{\linewidth}\raggedright
Candidate state \(X\)
\end{minipage} & \begin{minipage}[b]{\linewidth}\raggedright
\(\Delta\) bits/bit (held-out)
\end{minipage} & \begin{minipage}[b]{\linewidth}\raggedright
Verdict
\end{minipage} \\
\midrule\noalign{}
\endhead
\bottomrule\noalign{}
\endlastfoot
Wheel \(S_{n-1}\) & \(+0.000\) & No signal \\
\(a_{n-1} \bmod m\), any \(m \le 1024\) & \(\approx 0.000\) or negative
& No signal \\
\(n \bmod m\), any \(m\) & \(-0.032\) & Actively overfits \\
Joint \((a_{n-1} \bmod m,\, n \bmod m)\) & \(\le 0\) for \(m \le 256\) &
No signal \\
Visited-density bucket \(\rho_n\) & \(+0.000\) & No signal \\
\textbf{Previous bit \(b_{n-1}\)} & \(\mathbf{+0.956}\) &
\textbf{Near-complete} \\
Last 2 bits \((b_{n-2}, b_{n-1})\) & \(+0.968\) & Saturating \\
Last 3 bits & \(+0.971\) & Saturating \\
Last 16 bits & \(+0.972\) & No further lift \\
\end{longtable}

The obstruction stream is \textbf{autoregressive on its own history},
not on any positional or residue structure of the value sequence. The
direct joint table at \(N = 2 \times 10^5\) is:

\[
\begin{array}{c|cc}
& b_n = 0 & b_n = 1 \\ \hline
b_{n-1} = 0 & 720 & 99{,}269 \\
b_{n-1} = 1 & 99{,}270 & 740
\end{array}
\]

equivalently \[
\Pr[b_n \ne b_{n-1}] \;=\; 0.993, \qquad \Pr[b_n = b_{n-1}] \;=\; 0.007.
\]

\textbf{Recamán's obstruction stream is an almost-perfect alternating
sequence \(\cdots\, 0\, 1\, 0\, 1\, 0\, 1 \cdots\) with rare phase-slip
defects at rate \(\sim 7 \times 10^{-3}\).} The down/up balance is not a
stationary distribution over a non-trivial state space; it is the
symmetry of the alternation. The wheel's symbolic flip was modelling a
process whose actual dynamics are \emph{deterministic alternation plus
rare defects}.

\hypertarget{the-phase-slip-process}{%
\subsection{3.5 The phase-slip process}\label{the-phase-slip-process}}

What remains to model is the location of the phase-slip events. Define
\[
\sigma_n \;=\; b_n \oplus b_{n-1} \;\in\; \{0, 1\},
\] the indicator that the alternation \emph{broke} at step \(n\)
(\(\oplus\) is XOR). Then \(\sigma_n = 0\) means ``alternation continued
normally''; \(\sigma_n = 1\) means ``phase slipped.'' We have
\(\Pr[\sigma_n = 1] \approx 0.007\), and the problem of describing
Recamán's obstruction process \textbf{reduces to describing where these
rare events occur}.

The triple counts of \((b_{n-2}, b_{n-1}, b_n)\) at
\(N = 2 \times 10^5\) are illuminating:

\begin{longtable}[]{@{}lll@{}}
\toprule\noalign{}
Triple & Count & Fraction \\
\midrule\noalign{}
\endhead
\bottomrule\noalign{}
\endlastfoot
\((0,1,0)\) & \(98{,}937\) & \(0.4947\) \\
\((1,0,1)\) & \(98{,}761\) & \(0.4938\) \\
\((1,0,0) = (0,0,1)\) & \(508\) each & \(0.0025\) each \\
\((0,1,1) = (1,1,0)\) & \(\sim 333\) each & \(0.0017\) each \\
\((0,0,0), (1,1,1)\) & \(212, 407\) & \(0.0011, 0.0020\) \\
\end{longtable}

The two dominant triples \((0,1,0)\) and \((1,0,1)\) together account
for \(98.85\%\) of all triples. The defect triples are roughly twice as
common when crossing through \(b = 0\) as through \(b = 1\) --- a small
asymmetry that warrants further study.

\begin{quote}
\textbf{Restated open problem.} The closure problem of Section 6 is now
sharpened: characterise the \emph{phase-slip rate}
\(\Pr[\sigma_n = 1 \mid \text{state}]\) and find a state --- not the
\(\Theta_3\) wheel --- on which it is well-described. Candidate state
variables to test: the visited-set density at the candidate down-target
\(a_{n-1} - n\) (which the held-out scan did not yet test because we
used a global density proxy), the gap structure of \(V_{n-1}\) around
\(a_{n-1} - n\), or features of \(D_n\) from the signed-triangular
identity of Section 1.1.1.
\end{quote}

\hypertarget{layered-coordinates-and-the-markov-shadow-defect}{%
\section{4. Layered Coordinates and the Markov-Shadow
Defect}\label{layered-coordinates-and-the-markov-shadow-defect}}

\hypertarget{encoding}{%
\subsection{4.1 Encoding}\label{encoding}}

Fix a positional encoding \(E(h) = (d_0, \dots, d_{r-1})\) for candidate
values \(h\). \textbf{The encoding must be fixed before testing} to
avoid post-hoc tuning.

Define centered residues \[
u_p(h) \;=\; \min\!\bigl(h \bmod p,\; p - (h \bmod p)\bigr),
\] and shifted coordinates \[
X(h) \;=\; 1 + u_2(h), \qquad Y(h) \;=\; 1 + u_3(h).
\] The shift by \(1\) avoids the degenerate case \(XYZ = 0\) that
plagues the unshifted version.

A positional dipole score is \[
Z(h) \;=\; 1 + \sum_{0 \le i < j < r} \omega_{j-i}\,\mathbf{1}[d_i = d_j],
\] with fixed weights \(\omega_k \ge 0\).

\hypertarget{the-defect}{%
\subsection{4.2 The defect}\label{the-defect}}

The \textbf{Markov-shadow defect} is \[
\boxed{\;\Delta_\kappa(h) \;=\; X(h)^2 + Y(h)^2 + Z(h)^2 - \kappa\, X(h)\, Y(h)\, Z(h).\;}
\] Small \(|\Delta_\kappa(h)|\) means \(h\) lies near the surrogate
Markov surface --- by analogy with the integer Markov equation \[
x^2 + y^2 + z^2 \;=\; 3xyz
\] and its Vieta jump \(z \mapsto 3xy - z\).

\hypertarget{the-vieta-shadow-partner}{%
\subsection{4.3 The Vieta-shadow
partner}\label{the-vieta-shadow-partner}}

The (rounded) Vieta partner is \[
Z'_\kappa(h) \;=\; \bigl\lfloor \kappa\, X(h)\, Y(h) - Z(h) \bigr\rceil.
\] The \textbf{stronger Vieta-shadow test} asks whether \(Z'_\kappa(h)\)
matches the \(Z\)-coordinate of another recorded hole more often than
matched random controls predict.

\hypertarget{why-zeta3-cannot-be-the-exact-coupling}{%
\section{\texorpdfstring{5. Why \(\zeta(3)\) Cannot Be the Exact
Coupling}{5. Why \textbackslash zeta(3) Cannot Be the Exact Coupling}}\label{why-zeta3-cannot-be-the-exact-coupling}}

A natural cubic lift is \[
x^2 + y^2 + z^2 \;=\; \zeta(3)\, xyz.
\] If integers \(x, y, z \in \mathbb{Z} \setminus \{0\}\) satisfied
this, then \[
\zeta(3) \;=\; \frac{x^2 + y^2 + z^2}{xyz} \;\in\; \mathbb{Q},
\] contradicting \textbf{Apéry's theorem} that \(\zeta(3)\) is
irrational. Hence:

\begin{itemize}
\tightlist
\item
  No exact integer Markov equation with coupling \(\zeta(3)\) exists.
\item
  The \(\zeta(3)\)-surface is admissible only as an \textbf{approximate
  real surface} with a rounding defect, tracked via
  \(\lfloor \zeta(3)\, XY - Z \rceil\).
\end{itemize}

The empirical question is whether holes have smaller
\(|\Delta_{\zeta(3)}|\) than controls, \emph{not} whether they lie
exactly on the surface.

\hypertarget{the-closure-problem-reformulated}{%
\section{6. The Closure Problem
(Reformulated)}\label{the-closure-problem-reformulated}}

The closure problem \textbf{as originally posed} asked for
\(q(w) = \Pr[b_n = 1 \mid S_n = w]\) on the \(\Theta_3\) wheel. Section
3.3 showed this question is ill-posed:
\(q(0) \approx q(1) \approx 0.5\), so the wheel does not partition the
obstruction events. The closure variable is the wrong one.

The closure problem \textbf{as it should be posed}, given the empirical
finding of Section 3, is:

\begin{quote}
\textbf{Closure (phase-slip form).} Characterise the conditional
probability of a phase slip, \[
\pi_n(X) \;=\; \Pr\!\left[\sigma_n = 1 \,\middle|\, X_n\right],
\] where \(\sigma_n = b_n \oplus b_{n-1}\) is the slip indicator and
\(X_n\) is a candidate state variable computed from
\((a_{n-1}, n, V_{n-1})\).
\end{quote}

Marginally \(\Pr[\sigma_n = 1] \approx 0.007\). The closure problem is
to find a state \(X_n\) on which \(\pi_n(X_n)\) is concentrated on a
small subset and predictable on a held-out window.

The held-out scan of Section 3.4 ruled out the natural global
candidates: \(a_{n-1} \bmod m\), \(n \bmod m\), their joint, the wheel
state, and the global visited-set fill density. Each gave
\(\Delta \le 0\) bits/bit on held-out data. The candidates that remain
to be tested are \emph{local} features of \(V_{n-1}\):

\begin{enumerate}
\def\labelenumi{\arabic{enumi}.}
\tightlist
\item
  \textbf{Local gap structure.} The size and signed offset of the gap in
  \(V_{n-1}\) containing the candidate down-target \(a_{n-1} - n\).
\item
  \textbf{Neighbour density.}
  \(|V_{n-1} \cap [a_{n-1} - n - w,\; a_{n-1} - n + w]|\) for a small
  window \(w\).
\item
  \textbf{Last-collision recency.} The smallest \(k\) such that
  \(a_{n-1} - n + k \in V_{n-1}\) or \(a_{n-1} - n - k \in V_{n-1}\).
\item
  \textbf{Derivative of \(D_n\).} Whether the most recent down-step
  indices in \(D_n\) (from the signed-triangular identity, Section
  1.1.1) cluster near \(n\) or are spread.
\end{enumerate}

Each of these is local to the candidate down-target and therefore
plausibly predictive of whether the slip occurs at step \(n\). None has
been measured yet.

The scalar stationary density of \(a_n/n\) retains a role but as a
derived quantity. A flexible parametric family for held-out fitting is
\[
f_X(x) \;\propto\; \frac{1}{x^{\,p}\,(1 + x)^{\,q}}, \qquad p, q > 0.
\] Empirically, naive \(p = q = 1\) gives \(\mathbb{E}[X] \approx 0.70\)
on the support \([0.1, 4.04]\), while the histogram suggests
\(\mathbb{E}[X]\) closer to \(0.82\). The exponents must be fit on
held-out data, and the fit is downstream of --- not a substitute for ---
the phase-slip closure above.

\hypertarget{delannoy-structure-generated-by-recamuxe1n}{%
\section{7. Delannoy Structure Generated by
Recamán}\label{delannoy-structure-generated-by-recamuxe1n}}

\hypertarget{the-three-way-branch-is-intrinsic}{%
\subsection{7.1 The three-way branch is
intrinsic}\label{the-three-way-branch-is-intrinsic}}

At every step, Recamán's obstruction rule chooses among three local
outcomes: \[
\begin{aligned}
& \text{(D)}\quad \text{down-free:}\;\; a_{n-1} - n > 0 \text{ and } a_{n-1} - n \notin V_{n-1}, \\
& \text{(F)}\quad \text{forced-up by collision:}\;\; a_{n-1} - n > 0 \text{ but } \in V_{n-1}, \\
& \text{(B)}\quad \text{forced-up by boundary:}\;\; a_{n-1} - n \le 0.
\end{aligned}
\] The three outcomes commute pairwise in their local effect on the
operator state: (D) leaves the wheel fixed, (F) and (B) flip it via
\(\Theta_3\). The number of distinct operator-word trajectories of
length \(n\) over the alphabet \(\{\text{D}, \text{F}, \text{B}\}\),
subject to the wheel's two-state constraint, satisfies a three-term
recurrence \[
N(n) \;=\; \alpha N(n-1) + \beta N(n-2) + \gamma N(n-1, n-1),
\] which reduces to the Delannoy form \[
D(m, n) \;=\; D(m-1, n) + D(m, n-1) + D(m-1, n-1)
\] when the trajectory is indexed jointly by (D-count, blocked-count).

\hypertarget{what-this-means-for-the-positional-distribution}{%
\subsection{7.2 What this means for the positional
distribution}\label{what-this-means-for-the-positional-distribution}}

The Delannoy magnitudes appearing in the empirical positional
distribution of \(\{a_n\}\) are not an external envelope --- they are
\textbf{the trajectory count of Recamán's own obstruction process}. The
visited-set memory \(V_{n-1}\) enters as the \emph{measure} on these
trajectories: it determines which Delannoy paths carry mass and which
are nearly empty. Equivalently: \[
\#\{\text{Recamán trajectories of length } n\} \;\sim\; D(n, n) \cdot \mu_n,
\] where \(\mu_n\) is the path-mass measure induced by the memory term.

So the Delannoy structure is \textbf{intrinsic} to Recamán's rule; the
memory \(V_{n-1}\) reweights it but does not produce it. The ``drift
apart'' at large \(n\) reported in visualisations is the reweighting
becoming non-uniform --- not Delannoy ceasing to apply, but the measure
\(\mu_n\) concentrating on a small subset of the Delannoy paths.

\hypertarget{what-is-claimed}{%
\subsection{7.3 What is claimed}\label{what-is-claimed}}

\begin{quote}
\textbf{Claim (Delannoy generation).} The number of distinct
operator-word trajectories of length \(n\) generated by Recamán's
three-way obstruction branch satisfies a Delannoy recurrence. The
positional distribution of \(\{a_1, \dots, a_n\}\) inherits the
magnitude scaling \[
\log \#\{\text{trajectories}\} \;\sim\; n \log(1 + \sqrt{2})
\] with a memory-dependent reweighting measure \(\mu_n\) that becomes
non-uniform as \(|V_{n-1}|\) saturates.
\end{quote}

This is the structural statement that the Gemini visualisations were
tracking. The reweighting \(\mu_n\) is the actual object of interest at
large \(n\) --- it is what the closure problem (Section 6) must
determine.

\hypertarget{empirical-density-distribution-at-n-2-times-105}{%
\subsection{\texorpdfstring{7.4 Empirical density distribution at
\(N = 2 \times 10^5\)}{7.4 Empirical density distribution at N = 2 \textbackslash times 10\^{}5}}\label{empirical-density-distribution-at-n-2-times-105}}

We generated the full Recamán sequence to \(N = 200{,}000\) terms
(verifying the signed-triangular identity of Section 1.1.1 at every
checkpoint) and computed the density histogram with \(99\) bins,
matching the methodology of the screenshot. Summary statistics:

\begin{longtable}[]{@{}lll@{}}
\toprule\noalign{}
Statistic & \(N = 50{,}000\) & \(N = 200{,}000\) \\
\midrule\noalign{}
\endhead
\bottomrule\noalign{}
\endlastfoot
Max value \(a_n\) & \(201{,}976\) & \(1{,}345{,}032\) \\
Avg value & \(40{,}932\) & \(172{,}824\) \\
Peak density (mode bin) & \(7{,}183\) & \(29{,}797\) \\
Down-step fraction \(\|D_n\|/n\) & \(\approx 0.500\) & \(0.4995\) \\
\end{longtable}

\begin{figure}
\centering
\includegraphics[width=0.92\textwidth,height=\textheight]{recaman_compare.png}
\caption{Recamán density distribution at 50k and 200k terms. Top:
reproduces the screenshot exactly. Bottom: extended run with new
structure visible.}
\end{figure}

Three empirical observations:

\begin{enumerate}
\def\labelenumi{\arabic{enumi}.}
\tightlist
\item
  \textbf{The structural shape persists across \(N\).} Both panels show
  a sharp dominant peak, a secondary peak at roughly half the dominant's
  position, and a long sparse tail with smaller clusters. This is a
  \emph{fixed-shape distribution under rescaling}, consistent with the
  existence of a stationary law for \(a_n / n\) (Section 6).
\item
  \textbf{Peak positions shift with \(N\).} The dominant peak moves from
  \(\sim 62{,}000\) at \(N = 50{,}000\) to \(\sim 195{,}000\) at
  \(N = 200{,}000\) --- a factor of \(\approx 3.1\) for a factor-\(4\)
  increase in \(N\). This sublinear scaling is consistent with the
  heuristic \(a_n / n \to \text{const}\) but with persistent
  fluctuations.
\item
  \textbf{Secondary clusters near \(0.35 \cdot a_{\max}\)} appear at
  both scales and survive the extension. These are candidate ``fake
  plateau'' regions where the obstruction process produces extended runs
  of similar values.
\end{enumerate}

The down-step fraction \(|D_n|/n \approx 0.500\) at both scales is
structurally striking: it says the Recamán process is \textbf{balanced}
between up and down moves on the long run, even though the obstruction
rule is highly asymmetric in its instantaneous logic. The balance
presumably emerges from the visited-set density saturating at \(\rho^*\)
(the constant from Section 7.2).

\hypertarget{the-recamuxe1n-markov-shadow-conjecture}{%
\section{8. The Recamán Markov-Shadow
Conjecture}\label{the-recamuxe1n-markov-shadow-conjecture}}

\begin{quote}
\textbf{Conjecture (Recamán Markov-Shadow).} Let \(H\) be a set of late
or missing Recamán values and let \(C\) be a matched control set of
nearby or size-matched values that appear normally. For a fixed encoding
\(E\), fixed dipole weights \(\omega_k\), and fixed coupling \(\kappa\),
the values in \(H\) are statistically enriched for small Markov-shadow
defect \[
|\Delta_\kappa(h)| \;=\; \bigl|\,X(h)^2 + Y(h)^2 + Z(h)^2 - \kappa\, X(h)\, Y(h)\, Z(h)\,\bigr|
\] relative to \(C\).
\end{quote}

A stronger \textbf{Vieta-shadow test} asks whether \(Z'_\kappa(h)\)
matches the \(Z\)-score of another recorded hole more often than
expected under matched controls.

\hypertarget{empirical-protocol}{%
\section{9. Empirical Protocol}\label{empirical-protocol}}

The following protocol must be followed \emph{in order}, with no
parameter tuning on the test set:

\begin{enumerate}
\def\labelenumi{\arabic{enumi}.}
\tightlist
\item
  \textbf{Pre-register the encoding.} Fix \(E\), \(\omega_k\), and
  \(\kappa\) before looking at hole labels. Record them.
\item
  \textbf{Compute features.} For each \(h\), compute
  \(X(h), Y(h), Z(h), \Delta_\kappa(h), Z'_\kappa(h)\), and the wheel
  score \[
  W(h) \;=\; \operatorname{dist}\!\bigl(L(E(h)),\; \Theta_3(R(E(h)))\bigr).
  \]
\item
  \textbf{Match controls.} For each known hole, draw a size-matched
  control of integers of the same digit length not in the hole set.
\item
  \textbf{Test.} Compare distributions of \(|\Delta_\kappa|\), \(W\),
  and \(|Z'_\kappa - \text{nearest hole's }Z|\) for holes vs.~controls.
  Use permutation tests (10,000 reshuffles) and report two-sided
  \(p\)-values.
\item
  \textbf{Bonferroni correct} across the three statistics. Declare
  success only at \(p < 0.05/3\).
\item
  \textbf{Test the Delannoy generation.} Count distinct operator-word
  trajectories \(\#\{S_1, \dots, S_n\}\) produced by the true Recamán
  obstruction stream and verify the recurrence
  \(N(n) \sim D(\lfloor n/2 \rfloor, \lceil n/2 \rceil)\) at the
  magnitude level. Separately, estimate the reweighting measure
  \(\mu_n\) empirically and check whether its non-uniformity tracks
  \(|V_{n-1}| / \max V_{n-1}\).
\end{enumerate}

\begin{quote}
\textbf{Honest status.} As of this writing, steps 1--5 have not been
executed against the Chan hole list. A synthetic Markov simulation
exists but tests the wheel against itself, not real Recamán data. Step 6
is informally supported by visualisations showing Delannoy magnitudes
emerging from the Recamán positional distribution; the explicit
trajectory count and the reweighting measure \(\mu_n\) have not been
computed.
\end{quote}

\hypertarget{open-problems}{%
\section{10. Open Problems}\label{open-problems}}

\begin{enumerate}
\def\labelenumi{\arabic{enumi}.}
\tightlist
\item
  \textbf{Phase-slip closure.} Find a local state \(X_n\) (e.g.~local
  gap structure of \(V_{n-1}\) near \(a_{n-1} - n\)) on which
  \(\Pr[\sigma_n = 1 \mid X_n]\) is predictable on a held-out window.
  The global states (\(a_{n-1} \bmod m\), \(n \bmod m\), wheel) have
  already been falsified.
\item
  \textbf{Delannoy reweighting measure.} Determine the path-mass measure
  \(\mu_n\) on Recamán-generated Delannoy trajectories. Characterize how
  it concentrates as \(|V_{n-1}|\) saturates.
\item
  \textbf{Hole vs.~fake separation.} Decide whether \(930557\) is a real
  hole or a late fake using \(|\Delta_\kappa|\) and \(W\) as a
  classifier. Compare against the established hole \(852655\).
\item
  \textbf{Operator-word periodicity.} Is \(\{S_n\}\) eventually periodic
  for the true Recamán sequence?
\item
  \textbf{Stationary exponents.} Fit \(p, q\) in
  \(f_X(x) \propto x^{-p}(1+x)^{-q}\) on a held-out window and validate
  on the next.
\item
  \textbf{Markov-tree connection.} Does \(Z'_\kappa(h)\) --- the Vieta
  partner --- generate a tree of related holes, analogous to the Markov
  tree of integer triples?
\end{enumerate}

\hypertarget{honest-verdict}{%
\section{11. Honest Verdict}\label{honest-verdict}}

The synthesised drafts contained a mix of:

\begin{itemize}
\tightlist
\item
  A \textbf{genuine and underused taxonomy} (Section 2 ---
  local/global/free obstruction-bit classification).
\item
  A \textbf{legitimate empirical program} (Section 9 --- testing
  \(\Delta_\kappa\) and \(W\) against real holes with matched controls).
\item
  A \textbf{strong empirical negative result} (the \(\Theta_3\) wheel
  was falsified as a Markov model of the obstruction process by
  mutual-information, run-length, and KL-divergence tests at
  \(N = 2 \times 10^5\)).
\item
  A \textbf{clean positive empirical replacement} (the obstruction
  process is near-perfect alternation \(b_n \ne b_{n-1}\) with
  phase-slip rate \(\sim 7 \times 10^{-3}\), established by held-out
  validation).
\item
  A \textbf{structural identification} (Delannoy magnitudes are
  generated by Recamán's intrinsic three-way branch, reweighted by the
  memory measure \(\mu_n\)).
\item
  One \textbf{clean negative theoretical result} (Apéry blocks any exact
  \(\zeta(3)\)-Markov surface).
\end{itemize}

What does not survive: the wheel as a Markov model of the obstruction
process (falsified); softmax-as-attention (vacuous for two states); the
seven-approaches list (most aspirational); the earlier ``external
envelope'' framing of Delannoy (which inverted the direction of the
relationship); the original \texttt{racamman\_markov.py} script (its
block-probability generator is circular --- see Appendix B).

The framework lives or dies on two fronts now. \textbf{On the value
side}: until \(|\Delta_\kappa|\) or \(W\) actually separates holes from
controls on the Chan list (Section 9), the Markov-shadow conjecture is a
feature-engineering proposal. \textbf{On the obstruction-process side}:
until a local state \(X_n\) is found on which the phase-slip rate
\(\pi_n(X_n)\) is predictable on held-out data (Section 6), the closure
problem stands open. The taxonomy in Section 2 stands independently.

\hypertarget{appendix-a.-the-branch-map-identity-in-compact-form}{%
\section{Appendix A. The Branch-Map Identity in Compact
Form}\label{appendix-a.-the-branch-map-identity-in-compact-form}}

For all three systems: \[
\boxed{\; x_{n+1} \;=\; F_{b_n}(x_n), \qquad b_n \;=\; \Phi_n\bigl(\text{state},\,\text{history}\bigr). \;}
\] The function \(\Phi_n\) is:

\begin{itemize}
\tightlist
\item
  \(\Phi_n^{\text{Collatz}}(x) = x \bmod 2\) --- pointwise, memoryless.
\item
  \(\Phi_n^{\text{Recamán}}(a, V) = \mathbf{1}[a - n \le 0 \,\lor\, a - n \in V]\)
  --- non-local, memory-dependent.
\item
  \(\Phi_n^{\text{Delannoy}}\) --- external choice, unconstrained.
\end{itemize}

Hardness of the inverse problem (predict \(b_n\)) scales with the
information content of \(\Phi_n\): \[
H\bigl(\Phi^{\text{Delannoy}}\bigr) \;=\; \log 3, \qquad
H\bigl(\Phi^{\text{Collatz}}\bigr) \;=\; 1, \qquad
H\bigl(\Phi^{\text{Recamán}}\bigr) \;\le\; \log |V_{n-1}|.
\] Recamán's \(\Phi\) has unbounded information content; this is the
formal sense in which it is ``the hardest'' of the three.

\hypertarget{appendix-b.-reference-implementation-notes}{%
\section{Appendix B. Reference Implementation
Notes}\label{appendix-b.-reference-implementation-notes}}

Two scripts accompany this note. The first
(\texttt{racamman\_markov.py}, the original) is \textbf{kept for
archival reference only} --- its block-probability generator is a
hand-crafted Bernoulli rule, and the script trains on data the same rule
produced, which is circular and proves nothing. Worse, its assumption
that \(b_n\) can be drawn as a Bernoulli random variable contradicts
Recamán's defining property: \(b_n\) is \emph{deterministic} given the
full visited set \(V_{n-1}\).

The replacement scripts are:

\begin{itemize}
\item
  \textbf{\texttt{recaman\_wheel\_honest.py}}: runs the true Recamán
  generator with the full visited set, extracts the real obstruction
  stream, measures \(q(w) = \Pr[b_n = 1 \mid S_{n-1} = w]\) on the
  wheel, simulates a Markov null using those measured probabilities, and
  compares the two via run-length distributions and \(k\)-block KL
  divergence. Produces the numbers reported in Section 3.3.
\item
  \textbf{\texttt{recaman\_modm\_scan.py}} and
  \textbf{\texttt{recaman\_heldout.py}}: enumerate candidate state
  variables (\(a_{n-1} \bmod m\), \(n \bmod m\), joint, density proxy,
  bit history) and score each via mutual information and held-out
  log-likelihood. Produces the table in Section 3.4 and identifies
  bit-history conditioning as the only candidate that beats the marginal
  on held-out data.
\end{itemize}

Reproducibility: both scripts run deterministically on the true Recamán
sequence (no randomness in the generator). The Markov null comparison
uses a fixed seed. All numbers in Sections 3.3, 3.4, and 7.4 of this
note are reproducible by running the scripts as-is to
\(N = 2 \times 10^5\).

\begin{center}\rule{0.5\linewidth}{0.5pt}\end{center}

\emph{End of note.}

\end{document}
